\documentclass[12pt,a4paper]{article}
\usepackage[utf8]{inputenc}
\usepackage[T1]{fontenc}
\usepackage{amsmath}
\usepackage{color}
\usepackage{amsfonts}
\usepackage{amssymb}
\usepackage{graphicx}
\usepackage{cite}
\usepackage{geometry}
\usepackage{amsthm}
\geometry{margin=1in}


\newtheorem{assu}{Assumption}
\newtheorem{prop}{Proposition}

\title{Poisson directions.}
\date{\today}

\begin{document}

\section{Efficiency}

\begin{assu}
    Suppose $\mathbb{E}[Y] = \lambda(x, w)$, $\mathbb{V}[Y] = \sigma^2\lambda(x, w)$ is a sum of i.i.d. random variables coming from the same distribution.
\end{assu}


\begin{prop}
    Under this assumption, if we want to estimate $E[\phi_{\theta}(X)\tau] / E[\phi_{\theta}(X) Y(0)]$ or $E[\phi(X; P)\tau] / E[\phi(X; P)Y(0)]$, then optimizing for variance, this is the optimal choice for $w$.
\end{prop}

Think of putting $P$ to be the distribution of $Y(0)$ given $X$ instead?





What I do not like about this proposition is that why the hell do we allow w to depend on theta? I also wonder if it is possible to make the assumption a bit more abstract...


\section{Ex-post weights under model assumption}

\begin{proposition}
    Suppose $\mathbb{E}[Y(0) | X] = \exp(X\gamma)$, then the estimate of the Poisson regression is a weighted average in the following way:
\begin{equation}
    \theta = \frac{E[\psi_{1, \theta}(X, \theta)\tau]}{E[\psi_{0, \theta} Y(0)]}
\end{equation}
\end{proposition}

\begin{proof}
    From FOC with respect to $\beta$:
    \[\mathbb{E}[\exp(X\gamma)(1 + \theta_i W) - \exp(X\beta)(1 + \theta e(x))\mid X] = 0\]
    \[\exp(X\beta) = \mathbb{E}[\frac{\exp(X\gamma)(1 + \theta_i W)}{(1 + \theta e(x))} - \mid X] = 0\]


    Now plug in FOC w.r.t. $\theta$:

    \[\exp(X\beta) = \mathbb{E}[\frac{\exp(X\gamma)(1 + \theta_i W)}{(1 + \theta e(x))} - \mid X] = 0\]
\end{proof}

Done.

I need to assume exponent of the baseline
from X focs take exp(Xbeta), plug it into W foc
replace exp(Xgamma) with Y(0)
The 1 related part with cancel out (under the ex-ante weights) -- with somehting
The ex-post weights will remain.

I think I need tilted model based assumption

\section{Axiomatization}

It would be cool to motivate GLM as well. Suppose we want that the weights depend only on a finite statistic of Y? Can we get a Pitman-Koopman-Darmois type result that then we should have GLM?

Specify better, what do you mean here

\section{Empirical part}

\section{}

\end{document}
